\section{Introduction}
\label{sec:intro}

Robots that guide visitors, patrol or carry goods in public buildings share their space with
crowds that do not watch out for them~\cite{burgard1999museum,triebel2016spencer,
kruse2013humanaware}. Quadruped platforms are attractive there: they climb thresholds and
ramps and are perceived as approachable, but their bodies pitch and roll with the gait, their
footprint is long compared with its width, and their drivers expose only a velocity command.
Products built on such platforms tend to grow their autonomy around one robot: the footprint is
repeated in five parameter files, the base height is assumed in two different values, frames
are hard-coded, and safety checks that were right for one speed are silently wrong for the next.

Dharma grew out of such a product (a tour-guide robot) and was refactored into a general,
robot-agnostic stack. An independent review of the original code modelled every subsystem and
found three structural weaknesses: pose estimates carried no uncertainty and were never
corrected continuously against the map; geometric assumptions (robot radius, base height,
level frames) were inconsistent across packages; and the robot description was scattered, so
the stack was not robot-agnostic in practice. It also found safety defects: a planner radius
smaller than the robot's half diagonal, a protective field that ignored rotation, and blind
sensors read as free space. This paper describes the redesigned stack (v0.2.0) and the models
behind it.

\textbf{Contributions.}
\begin{itemize}
  \item A layered, robot-agnostic architecture in which one robot profile is the single source
    of geometry, limits and frames, with typed configuration validated at start and on every
    change, typed status interfaces and plugin registries (Sec.~\ref{sec:arch}).
  \item A gravity-aligned bridge for LiDAR-inertial odometry and a fiducial-anchored SE(2)
    extended Kalman filter over the map-to-odometry correction: drift grows with motion, not
    time, the innovation is taken where the robot is, and a chi-square gate opens with the
    uncertainty (Sec.~\ref{sec:loc}).
  \item One braking model shared by the speed profile, the path shift and a protective field
    that limits translation, rotation and combined motion, with fail-safe sensing semantics and
    a latched e-stop that must be re-armed (Secs.~\ref{sec:track}--\ref{sec:field}).
  \item The engineering that keeps the models consistent: pure algorithm cores, enforced module
    size and purity, a ROS~2 Humble and Jazzy CI matrix, and traceability of every review
    finding to a fix and a regression test (Sec.~\ref{sec:eng}).
  \item A seeded simulation protocol with odometry drift, latency and sensor dropout to evaluate
    stopping, blind-sensor behaviour, localization and avoidance (Sec.~\ref{sec:eval}).
\end{itemize}
